% Generated draft: requires mathematical and source review.
\documentclass[11pt]{article}
\usepackage[margin=25mm]{geometry}
\usepackage{fontspec}
\setmainfont{Latin Modern Roman}
\usepackage{amsmath,amssymb,mathtools,unicode-math}
\setmathfont{Latin Modern Math}
\usepackage{enumitem,needspace,xurl,hyperref}
\setlength{\parindent}{0pt}
\setlength{\parskip}{5pt}
\setlength{\emergencystretch}{3em}
\newcommand{\N}{\mathbb{N}}
\newcommand{\Z}{\mathbb{Z}}
\newcommand{\Q}{\mathbb{Q}}
\newcommand{\R}{\mathbb{R}}
\newcommand{\C}{\mathbb{C}}
\newcommand{\F}{\mathbb{F}}
\newcommand{\A}{\mathbb{A}}
\newcommand{\PP}{\mathbb{P}}
\newcommand{\E}{\mathbb{E}}
\newcommand{\eps}{\varepsilon}
\newcommand{\dd}{\,\mathrm{d}}
\begin{document}
\section*{UnsolvedMath ID: 30003673}
\textbf{Source problem code:} catalogue entry
\subsection{Definitions and mathematical statement}
The supplied reconstruction asks two questions for the fair-competition aggregation--diffusion model in spatial dimension \(N\in\N\), \(N>1\). It distinguishes a critical \emph{unrescaled} minimization problem from a subcritical \emph{rescaled} problem. The source record does not specify the functionals, the interaction kernel, the diffusion exponent, the admissible class, or the stationary equation. Consequently, the following is a conditional formulation with explicit placeholders, rather than a fully specified variational problem.

For a functional \(H\) on an admissible class \(\mathcal A\), write
\[
\operatorname{Argmin}_{\mathcal A}H
 =\{u\in\mathcal A:H(u)=\inf_{v\in\mathcal A}H(v)\}.
\]
Let \(\mathcal F_{\chi_c}\) denote the intended critical unrescaled free energy, \(\mathcal A_c\) its intended admissible class, and \(D_\lambda\), \(\lambda>0\), the intended dilation action. These objects, including the normalization and any centering restriction, remain to be supplied. Conditional on existence of a global minimizer, the critical question is whether
\[
\forall u,v\in\operatorname{Argmin}_{\mathcal A_c}\mathcal F_{\chi_c},
\qquad \exists\lambda>0\quad v=D_\lambda u.
\]
No existence assertion is established by this formulation.

For each interaction strength \(\chi\) in the intended subcritical range, let \(\mathcal G_\chi\) denote the rescaled functional, \(\mathcal A_{\mathrm r,\chi}\) its admissible class, and \(\mathcal S_\chi\subseteq\mathcal A_{\mathrm r,\chi}\) the stationary states of the corresponding rescaled equation. The subcritical range, rescaling, equation, and meaning of stationarity are not supplied. The additional questions in the reconstructed statement are whether
\[
\mathcal S_\chi\subseteq
\operatorname{Argmin}_{\mathcal A_{\mathrm r,\chi}}\mathcal G_\chi
\]
and whether, conditional on existence,
\[
\forall u,v\in
\operatorname{Argmin}_{\mathcal A_{\mathrm r,\chi}}\mathcal G_\chi,
\qquad u=v.
\]
The reconstructed subcritical question specifies uniqueness without a symmetry qualification; any intended normalization or quotient must be checked.

The original statement instead refers to a ``critical rescaled functional'' and asks only about uniqueness modulo dilations. This conflicts with the critical unrescaled formulation in the supplied reconstruction. The record reports that its reconstruction was checked against the official report, but that verification is supplied provenance and has not been independently performed here. The supplied dated assessment also identifies Newtonian and one-dimensional uniqueness results; thus its higher-dimensional open label must not be read as asserting that every higher-dimensional parameter case is unresolved.
\subsection{Short English statement}
For the intended fair-competition model in dimension \(N>1\), ask whether critical unrescaled global minimizers are unique up to dilation, and whether subcritical rescaled stationary states are global minimizers with a unique minimizer. The precise model, admissibility conditions, symmetry conventions, and a discrepancy concerning rescaling require clarification from the references.
\subsection{Sources}
\begin{itemize}
\item Supplied local UnsolvedMath registry record, permanent ID \(30003673\), exact source problem code \texttt{catalogue entry}, title ``Uniqueness of Higher-Dimensional Critical Minimizers.'' The record labels the problem open and its reconstructed statement \texttt{corrected\_verified}; these are supplied curation labels. No individual proposer is identified; the proposed year is recorded as 2017.
\item Supplied report citation: \emph{Classical and Quantum Mechanical Models of Many-Particle Systems} (2018), Oberwolfach Reports, pp.\ 3345--3425, DOI \texttt{10.4171/owr/2017/56}. \url{https://doi.org/10.4171/owr/2017/56}. The supplied background attributes the problem to workshop open-problem discussions documented in Oberwolfach Reports.
\item Supplied literature source: \emph{Equilibria in the fair-competition regime}. \url{https://arxiv.org/abs/1610.00939}.
\item Supplied literature source: \emph{Uniqueness of stationary Keller--Segel states}. \url{https://arxiv.org/abs/1905.07788}.
\item Supplied literature source, labeled ``OWR Report 56/2018'': \url{https://oa.tib.eu/renate/bitstreams/63dd963a-2f59-4c50-b626-ffc98e104dc8/download}.
\item Literature assessment dated exactly \texttt{2026-08-21}: the supplied curation states that uniqueness modulo dilations is proved in the Newtonian and one-dimensional cases and remains open in the higher-dimensional fair-competition regime. This is supplied dated curation derived from the OWR record, not an independently verified literature assessment; its time-sensitive status and cited claims require independent checking.
\end{itemize}
\end{document}
